\documentclass[12pt]{article}

\usepackage{amsmath}
\usepackage{amssymb}

\title{Tehtävät 2}

\begin{document}
\maketitle
\begin{enumerate}
\item[1.]
\begin{enumerate}
\item[(a)] Koska $Y_i \sim \text{Exp}(1/\mu)$, niin tämän tiheysfunktio on 
$$
f_{Y_i}(y_i; \mu) = \frac{1}{\mu} \exp(-\frac{1}{\mu} y_i),
$$
kun $y_i > 0$. Koska havainnot ovar riippumattomia, niin tilastollisen mallin lauseke eli yhteistiheysfunktio saadaan tulona
$$
\begin{aligned}
f(\textbf{y}; \mu) &= \prod_{i = 1}^n f_{Y_i}(y_i; \mu) \\
    &= \prod_{i = 1}^n \frac{1}{n} \exp(-\frac{1}{\mu} y_i) \\
    &= \frac{1}{\mu^n} \exp\Bigg( \sum_{i = 1}^n y_i \Bigg) \\
    &= \mu^{-n} \exp\Bigg( -\frac{1}{\mu}n\bar{y} \Bigg),
\end{aligned}
$$
kun $y_i > 0$ kaikilla $i$, missä 
$$
\bar{y} = \frac{1}{n} \sum_{i = 1}^n y_i.
$$
\item[(b)] Yhteistiheysfunktiosta saadaan suoraan aineistoa $\textbf{y}$ vastaava uskottavuusfunktio
$$
L(\mu; \textbf{y}) = \mu^{-1} \exp\Bigg( -\frac{1}{\mu} n \bar{y} \Bigg).
$$
Log-uskottavuusfunktio on tällöin
$$
\begin{aligned}
\ell(\mu; \textbf{y}) &= \log(L(\mu; \textbf{y}) \\
    &= \log \Bigg( \mu^{-n} \exp \bigg( -\frac{1}{\mu} n \bar{y} \bigg) \Bigg) \\
    &= -n \log(\mu) - \frac{1}{\mu} n\bar{y}.
\end{aligned}
$$

\end{enumerate}
\item[2.] 
\begin{enumerate}
\item[(a)] Log-uskottavuusfunktion derivaatta on
$$
\ell^\prime(\mu) = -n \frac{1}{\mu} + \frac{1}{\mu^2} n\bar{y} = \frac{n}{\mu^2}(-\mu + \bar{y}).
$$
Huomataan, että
$$
\begin{array}{lll}
& \ell^\prime &= 0 \\
\Leftrightarrow & -\mu + \bar{y} &= 0 \\
\Leftrightarrow & \mu &= \bar{y}.
\end{array}
$$
Koska lisäksi huomataan, että
$$
\ell^\prime(\mu) = \frac{n}{\mu^2}(-\mu + \bar{y}) > \frac{n}{\mu^2}(-\bar{y} + \bar{y}) = 0
$$
kaikilla $\mu < \bar{y}$ ja
$$
\ell^\prime(\mu) = \frac{n}{\mu^2}(-\mu + \bar{y}) < \frac{n}{\mu^2}(-\bar{y} + \bar{y}) = 0
$$
kaikilla $\mu > \bar{y}$, niin $\bar{y}$ on funktion $\ell$ (ja siten myös uskottavuusfunktion $L$) maksimikohta.
Siis su-estimaatti on
$$
\hat{\mu} = \bar{y}
$$
ja su-estimaattori on 
$$
\hat{\mu}(\textbf{Y}) = \bar{Y}.
$$
\item[(b)] Aineistoa $\textbf{y} = (2,25, 8,40, 5,10, 11,85, 3,30, 8,10)$ vastaavaksi su-estimaatiksi saadaan
$$
\hat{\mu} = \bar{y} = \frac{1}{6}\sum_{i = 1}^6 y_i = \frac{1}{6}(2,25 + 8,40 + 5,10 + 11,85 + 3,30 + 8,10) = 6,5.
$$
Jos $Y_i \sim \text{Exp}(1/\mu)$, niin $\mathbb{E}[Y_i] = \mu$. Täten saadun su-estimaatin perusteella voidaan päätellä, että koiraihmisten väliajan noudattaman eksponentijakauman odotusarvo on luultavasti suunnilleen suuruusluokkaa 6,5. Aineston perusteella 6,5 on uskottavin su-estimaatin arvo, joten 3 ei ole kovin uskottava, mutta se on kuitenkin arvo.
\end{enumerate}
\item[3.] 
\begin{enumerate}
\item[(a)] Nyt odotusarvon lineaarisuuden nojalla su-estimaattorin $\hat{\mu}(\textbf{Y}) = \bar{Y}$ odotusarvoksi saadaan
$$
\mathbb{E}_\mu[\hat{\mu}(\textbf{Y})] = \mathbb{E}[\bar{Y}] = \mathbb{E}_\mu[\frac{1}{n} \Bigg( \sum_{i = 1}^n Y_i \Bigg) = \frac{1}{n} \sum_{i = 1}^n \mathbb{E}_\mu[Y_i] = \frac{1}{n} \sum_{i = 1}^n \mu = \frac{1}{n}n\mu = \mu.
$$
Harhattomuuden Määritelmän 3.4 nojalla nähdään, että SU-estimaattori on harhaton, sillä määritelmä sanoo, ettää parametring $\theta$ estimaattori $\hat{\theta}(\textbf{Y})$ on harhaton \textit{(engl. biased)}, mikäli estimaattorin odotusarvo on sama kuin parametrin todellinen arvo, eli
$$
\mathbb{E}_\mu[\hat{\theta}(\textbf{Y})] = \theta, \text{ kaikilla } \theta.
$$
\item[(b)] Varianssin laskemisessa on käytetty tietoa lineaarimuunnoksen varianssista sekä riippumattomien satunnaismuuttujien $Y_1, \ldots, Y_n$ summan varianssista. Su-estimaattorin $\hat{\mu}(\textbf{Y}) = \bar{Y}$ varianssiksi saadaan siten
$$
\begin{aligned}
\text{Var}_\mu\hat{\mu}(\textbf{Y}) &= \text{Var}_\mu(\bar{Y}) \\
    &= \text{Var}_\mu \Bigg( \frac{1}{n} \sum_{i = 1}^n Y_i \Bigg) \\
    &= \frac{1}{n^2} \text{Var}_\mu\Bigg( \sum_{i=1}^n Y_i \Bigg) \\
    &= \frac{1}{n^2} \sum_{i = 1}^n \text{Var}_\mu(Y_i) \\
    &= \frac{1}{n^2}n\mu^2 \\
    &= \frac{1}{n} \mu^2. 
\end{aligned}
$$
\end{enumerate}
\item[4.] c
\item[5.] d
\end{enumerate}
\end{document}